% ============================================================
% Capitolo 6 — Fasi e Modes
% ============================================================

\chapter{Fasi e Modes}
\label{ch:fasi-e-modes}

Uno degli errori più comuni nell'uso degli agenti AI è trattarli come strumenti uniformi: la stessa configurazione per qualsiasi task, in qualsiasi momento del progetto.

Le \textbf{Fasi} dicono all'agente in che punto del ciclo di sviluppo si trova. I \textbf{Modes} controllano quanta cerimonia applicare per quella sessione. Insieme, rendono l'agente contestuale invece che generico.

\section{Le sette fasi del ciclo AI-Driven}
\label{sec:06-sette-fasi}

\begin{table}[htbp]
  \centering
  \begin{tabular}{cll}
    \toprule
    \textbf{Fase} & \textbf{Nome} & \textbf{Quando ci sei} \\
    \midrule
    0 & Discovery & Raccolta requisiti, comprensione del problema \\
    1 & Analysis & Analisi requisiti, scelta dell'approccio macro \\
    2 & Design & Architettura, API contract, schema DB \\
    3 & Implementation & Codice, test \\
    4 & Verification & QA, integration testing \\
    5 & Release & Deploy, documentazione finale \\
    6 & Ops & Monitoraggio, incident response, manutenzione \\
    \bottomrule
  \end{tabular}
  \caption{Le sette fasi del ciclo AI-DLC}
  \label{tab:06-fasi}
\end{table}

Il campo \texttt{Phase} in \texttt{\_CONTEXT.md} indica la fase corrente e determina quale modulo il framework carica:

\begin{table}[htbp]
  \centering
  \begin{tabular}{cll}
    \toprule
    \textbf{Fase} & \textbf{Modulo} & \textbf{Skill suggerite} \\
    \midrule
    0--1 & \texttt{02\_DISCOVERY\_ANALYSIS.md} & \texttt{SKILL\_ANALYSIS.md} \\
    2 & \texttt{03\_DESIGN.md} & \texttt{SKILL\_DESIGN.md}, \texttt{SKILL\_API\_DESIGN.md} \\
    3 & \texttt{04\_IMPLEMENTATION.md} & \texttt{SKILL\_API\_DESIGN.md}, \texttt{SKILL\_SECURITY.md} \\
    4--5 & \texttt{05\_VERIFICATION\_RELEASE.md} & \texttt{SKILL\_TESTING.md} \\
    6 & \texttt{06\_OPS.md} & \texttt{SKILL\_OPS.md} \\
    \bottomrule
  \end{tabular}
  \caption{Moduli e skill per fase}
  \label{tab:06-moduli-fasi}
\end{table}

\section{I cinque Modes}
\label{sec:06-cinque-modes}

\subsection*{LITE --- lavoro quotidiano}

Per lavoro routine su un progetto stabile. Riduce l'overhead senza rinunciare alla sicurezza.

\begin{itemize}
  \item Conferma di sessione: saltata se \texttt{\_CONTEXT.md} non è cambiato
  \item Confidence tag: solo per output $>$10 righe o classificati HIGH
  \item Checkpoint: suggerito ogni 5+ azioni (non obbligatorio per LOW/MEDIUM)
  \item HALT trigger: sempre rispettati
\end{itemize}

\subsection*{STANDARD --- default}

Applica il protocollo completo. Per nuovi progetti, nuovi team, zone del codice mai toccate.

\subsection*{AUDIT --- tracciabilità completa}

Per lavoro con compliance, sicurezza, revisioni formali. Confidence tag su ogni output, checklist di completamento per ogni task.

\subsection*{RAPID --- spike/emergenze}

Lavoro time-boxed (1--2 ore), nessuna documentazione obbligatoria. Usa sempre un branch separato; evita i path di produzione. SEC critici sempre rispettati.

\subsection*{FAST --- task semplici}

Output diretto senza loop di conferma, a meno che il rischio non sia HIGH.

\section{Scenario tipico su TaskFlow API}
\label{sec:06-scenario}

\begin{table}[htbp]
  \centering
  \begin{tabular}{lll}
    \toprule
    \textbf{Periodo} & \textbf{Mode} & \textbf{Motivazione} \\
    \midrule
    Fase 0--2 & STANDARD & Decisioni architetturali ad alto impatto \\
    Prima settimana Fase 3 & STANDARD & Nuovo territorio, stabilire convenzioni \\
    Settimane 2--N Fase 3 & LITE & Pattern rodati, team allineato \\
    Audit pre-release & AUDIT & Tracciabilità completa richiesta \\
    Hotfix urgente & RAPID & Velocità, branch dedicato \\
    Refactor di naming & FAST & Task a basso rischio \\
    \bottomrule
  \end{tabular}
  \caption{Evoluzione del Mode durante il progetto}
  \label{tab:06-scenario}
\end{table}

Per cambiare mode aggiorna il campo \texttt{Mode:} in \texttt{\_CONTEXT.md}. Esempio concreto:

Lorenzo, dopo due settimane in STANDARD, passa a LITE:

\begin{lstlisting}[language=,caption={Aggiornamento del Mode nel context}]
| Mode | LITE |
\end{lstlisting}

Alla sessione successiva, l'agente non chiede la conferma di bootstrap --- la presuppone invariata. Lorenzo scrive \textit{``Implementa \texttt{DELETE /tasks/:id}''} e l'agente propone subito il piano senza il passaggio di conferma contesto. Il flusso è più fluido, le regole di sicurezza sono le stesse.

\section{Errori comuni con Fasi e Modes}
\label{sec:06-errori}

\textbf{Tenere STANDARD per tutta la vita del progetto.} Dopo le prime settimane, le conferme diventano rumore. Il team inizia a ignorarle. Abbassa a LITE quando il progetto è stabile.

\textbf{Usare RAPID per lavoro che va in produzione.} RAPID è per esperimenti. Se il codice prodotto in RAPID finisce in main senza review, stai usando il mode sbagliato.

\textbf{Non aggiornare la fase in \texttt{\_CONTEXT.md} quando si cambia fase.} L'agente carica i moduli sbagliati. Il cambio di fase nel PROGRESS.md deve essere accompagnato dall'aggiornamento del campo \texttt{Phase} nel context.

\textbf{Pensare che i Modes disattivino gli HALT trigger.} Non è così. In nessun mode gli HALT trigger vengono disattivati.

\section*{\LabelSummary}

\begin{itemize}
  \item Le \textbf{sette fasi} (0--6) definiscono il punto del ciclo di sviluppo e quale modulo viene caricato.
  \item I \textbf{cinque Modes} (LITE, STANDARD, AUDIT, RAPID, FAST) controllano quanta cerimonia applica l'agente.
  \item Mode e Fase sono indipendenti: puoi essere in Fase 3 con Mode LITE.
  \item Gli HALT trigger non vengono mai disattivati, indipendentemente dal Mode.
  \item Cambia Mode liberamente durante il progetto: STANDARD per partire, LITE per il lavoro quotidiano rodato.
\end{itemize}
